\section{Research Methodology and Experimental Evaluation}

\towrite[Aim: open the chapter in one short paragraph. Cover: how the evaluation will answer the research question and test any hypotheses; the order of the sections that follow.]

\subsection{Research Design}

\towrite[Aim: name and justify the design. Cover: quantitative experimental evaluation, or a mixed approach if you add questionnaires; within-subject or repeated measures as applicable; independent and dependent variables; why this design fits the question.]

\subsection{Experimental Setup and Equipment}

\towrite[Aim: let another researcher rebuild the setup. Cover: headset and sensor preparation; the car; the computer and software versions; the physical arrangement, with a photo or diagram; the procedure for checking signal contact quality.]

\subsection{Car Construction and Electronics Integration}

\towrite[Aim: record the car as used in the experiments. Cover: the final configuration, with a photo; any changes made after Chapter 3; how consistent behaviour was checked before each session, such as battery level and motor response. Refer to Section 3.5 for the design rather than repeating it.]

\subsection{Participants and Ethical Considerations}

\towrite[Aim: describe who took part and how they were protected. Cover: recruitment; inclusion and exclusion criteria; number of participants and why; informed consent, with forms in Appendix B; ethics approval from the institution; data protection under GDPR; safety and the right to withdraw. If approval is pending, say so and do not state it as granted.]

\subsection{Experimental Environment}

\towrite[Aim: describe where the tests took place. Cover: the room or track; layout and dimensions; lighting, noise and sources of electrical interference; how conditions were kept consistent between sessions.]

\subsection{Command Training and Calibration}

\towrite[Aim: explain how participants prepared. Cover: which commands were trained; the training tool; number and length of training rounds; the criterion for being ready; how the same procedure was kept for everyone.]

\subsection{Experimental Tasks and Protocol}

\towrite[Aim: give the step-by-step procedure. Cover: each task, such as cued single commands and driving a set course; trial numbers and order; cue timing; rest breaks; rules for restarting or aborting. Put the full protocol in Appendix A.]

\subsection{Pilot Study}

\towrite[Aim: explain how the method was tested first. Cover: who took part; what you tried; what you changed as a result, such as thresholds, task length or instructions. Keep pilot data separate from the main results.]

\subsection{Main Experimental Sessions}

\towrite[Aim: describe how the main study ran. Cover: number of sessions and spacing between them; session structure and timing; who was present; how sessions were recorded; any deviations from the protocol.]

\subsection{Data Collected and Availability}

\towrite[Aim: state what data exists and who can access it. Cover: raw EEG if accessible, recognised commands, timestamps and the car's responses to tasks; file organisation and naming; pseudonymisation; secure storage and retention; what can be shared openly, with reproduction details in Appendix F. Check whether your Emotiv licence gives access to raw EEG.]

\subsection{Evaluation Metrics}

\towrite[Aim: define how performance is measured. Cover: how the metrics link to the hypotheses and objectives; shared definitions of trial, command window and correct response. For each metric below give a formula, the unit and the logged data used.]

\subsubsection{Command Success Rate}

\towrite[Aim: define this metric. Cover: correct commands executed as a share of commands cued; what counts as correct; the time window allowed.]

\subsubsection{False Activation Rate}

\towrite[Aim: define this metric. Cover: commands detected when none was intended, or a different one was; whether it is expressed per minute or as a share of trials, and why.]

\subsubsection{Missed Command Rate}

\towrite[Aim: define this metric. Cover: cued commands not detected within the time window, as a share of commands cued; how it relates to the success rate.]

\subsubsection{Command Detection Latency}

\towrite[Aim: define this metric. Cover: time from cue to detection by the framework; where each timestamp comes from; clock synchronisation; timing accuracy and how you verified it.]

\subsubsection{End-to-End Response Time}

\towrite[Aim: define this metric. Cover: time from cue to physical car response; how car movement is timestamped, such as onboard feedback or video; the breakdown into detection, processing and transmission delay.]

\subsubsection{Task Completion Rate}

\towrite[Aim: define this metric. Cover: tasks finished as a share of attempted tasks; the criteria for success and failure, such as time limits or leaving the course; any completion time measures.]

\subsubsection{Performance Variation Across Sessions}

\towrite[Aim: define how change is measured. Cover: the metrics compared across sessions; how learning, fatigue and variation between participants are examined; how you handle missing sessions.]

\subsection{Data Analysis Methods}

\towrite[Aim: explain the analysis before it is run. Cover: data cleaning and exclusion rules; descriptive statistics; statistical tests linked to each hypothesis, with assumption checks; significance level; effect sizes; software used. Choose tests suited to your sample size.]

\subsection{Optional Investigation of a Custom Machine Learning Approach}

\towrite[Aim: describe the optional comparison, if it is carried out. Cover: why a custom classifier could improve on built-in detection; available data; preprocessing and features; candidate models; validation that avoids leakage between sessions; comparison with Emotiv's detection. If it is not done, say so and remove Section 5.9.]

\subsection{Validity, Reliability, and Reproducibility}

\towrite[Aim: assess the method's strengths and threats. Cover: internal validity, such as learning and fatigue effects; external validity, such as sample and setting; reliability of measurement; how the open code, protocol and data allow reproduction.]

\subsection{Chapter Summary}

\towrite[Aim: close the chapter. Cover: the design, metrics and analysis in brief, leading into the results.]
